% =============================================================
% บทที่ 1 - บทนำ
% =============================================================

% -----------------------------------------------------------
\section{ความเป็นมาและความสำคัญของปัญหา}
% -----------------------------------------------------------

ในปัจจุบัน ข่าวสารเป็นข้อมูลสำคัญที่ช่วยให้ประชาชนรับรู้และติดตามเหตุการณ์ได้อย่างทันท่วงที โดยเฉพาะอย่างยิ่งในยุคที่ข้อมูลมีการเปลี่ยนแปลงตลอดเวลา อย่างไรก็ตาม ผู้อ่านที่ต้องการติดตามข่าวสารจากหลายแหล่งมักประสบปัญหาหลายประการ ได้แก่ ต้องเปิดเว็บไซต์หรือแอปพลิเคชันหลายแห่งเพื่อเปรียบเทียบข่าวเดียวกัน ทำให้เสียเวลาและเกิดความยุ่งยากในการติดตามสถานการณ์ที่เปลี่ยนแปลงอย่างรวดเร็ว นอกจากนี้ ข่าวแต่ละฉบับยังมีเนื้อหายาวและมีข้อมูลจำนวนมาก ทำให้ผู้อ่านต้องใช้เวลาในการอ่านเพื่อสกัดประเด็นสำคัญ รวมถึงยังต้องอ่านข่าวจากแหล่งภาษาอังกฤษเพื่อให้ได้มุมมองที่หลากหลาย อีกทั้งข่าวสารที่มีปริมาณมากและหลากหลายหมวดหมู่ยังทำให้การคัดกรองและวิเคราะห์ข่าวเป็นเรื่องท้าทาย

การนำเทคโนโลยีเว็บสแครปปิง (Web Scraping) และ RSS Feed มาใช้ในการดึงข้อมูลข่าวจากเว็บไซต์ข่าวอัตโนมัติ ตลอดจนการใช้เทคนิคการประมวลผลภาษาธรรมชาติ (Natural Language Processing: NLP) ในการจำแนกหมวดหมู่ข่าวตามประเภท และการใช้โมเดลภาษาขนาดใหญ่ (Large Language Model: LLM) ในการสรุปเนื้อหา จะช่วยลดภาระของผู้อ่านในการเปิดหลายเว็บไซต์และอ่านบทความฉบับเต็ม อีกทั้งยังสามารถนำเสนอข้อมูลสำคัญที่คัดสรรและวิเคราะห์แล้วแบบเรียลไทม์ได้ในที่เดียว

จากปัญหาดังกล่าว โครงงานนี้จึงมีแนวคิดในการพัฒนาระบบจำแนกและวิเคราะห์ข่าวสารอัตโนมัติด้วยเทคนิคการประมวลผลภาษาธรรมชาติแบบเฉพาะประเภท (Category-Specific NLP Analytics Pipeline) โดยระบบจะทำการดึงข่าวจากแหล่งข่าวเป็นระยะแบบอัตโนมัติ จำแนกหมวดหมู่ข่าวด้วยเทคนิค NLP และโมเดลแมชชีนเลิร์นนิง สรุปเนื้อหาด้วยโมเดลภาษาขนาดใหญ่ และแสดงผลข่าวพร้อมการวิเคราะห์บนเว็บแอปพลิเคชันแบบเรียลไทม์ผ่าน WebSocket เพื่อให้ผู้ใช้สามารถติดตามข่าวสารได้อย่างสะดวก รวดเร็ว และเหมาะสมกับความสนใจของแต่ละบุคคล

% -----------------------------------------------------------
\section{วัตถุประสงค์}
% -----------------------------------------------------------

\begin{enumerate}[label=\arabic*)]
    \item เพื่อศึกษาและวิเคราะห์เทคนิคการดึงข้อมูลข่าวจากเว็บไซต์ข่าวด้วยวิธีเว็บสแครปปิงและการอ่าน RSS Feed รวมถึงเทคนิคการประมวลผลภาษาธรรมชาติสำหรับการจำแนกหมวดหมู่ข่าวตามประเภท
    \item เพื่อออกแบบและพัฒนาระบบจำแนกและวิเคราะห์ข่าวสารอัตโนมัติด้วยเทคนิคการประมวลผลภาษาธรรมชาติแบบเฉพาะประเภท (Category-Specific NLP Analytics Pipeline) ครอบคลุมการดึงข่าว การจำแนกหมวดหมู่ การสรุปเนื้อหาด้วยโมเดลภาษาขนาดใหญ่ อัลกอริทึมแนะนำข่าว และการแสดงผลแบบเรียลไทม์
    \item เพื่อทดสอบและประเมินประสิทธิภาพของระบบในด้านความถูกต้องของการจำแนกหมวดหมู่ข่าว ความถูกต้องของการสรุปข่าว ความแม่นยำของอัลกอริทึมแนะนำข่าว และความสามารถในการแสดงผลข่าวแบบเรียลไทม์
\end{enumerate}

% -----------------------------------------------------------
\section{ขอบเขตและข้อจำกัดของระบบ}
% -----------------------------------------------------------

\subsection{ขอบเขต}

ระบบเป็นเว็บแอปพลิเคชันแบบเต็มรูปแบบ (Full-stack) ที่มีขอบเขตการทำงาน ดังนี้

\begin{enumerate}[label=\arabic*)]
    \item \textbf{การรวบรวมข่าวอัตโนมัติ:} ระบบสามารถดึงข่าวจากแหล่งข่าวออนไลน์หลายแห่ง โดยใช้เทคนิคเว็บสแครปปิง การอ่าน RSS Feed และการจำลองเบราว์เซอร์ด้วย Playwright สำหรับเว็บไซต์ที่ใช้ JavaScript
    \item \textbf{การสรุปข่าว:} ระบบส่งเนื้อหาข่าวในรูปแบบ Markdown ให้โมเดลภาษาขนาดใหญ่ผ่าน API ของมหาวิทยาลัยขอนแก่น (KKU Gen.AI) เพื่อสร้างบทสรุป ประเด็นสำคัญ อารมณ์ของข่าว (Sentiment) และคำสำคัญ โดยสรุปเป็นภาษาเดียวกับต้นฉบับข่าว
    \item \textbf{การจำแนกหมวดหมู่:} ระบบจำแนกข่าวออกเป็นหมวดหมู่เฉพาะทาง โดยใช้การตัดคำภาษาไทยด้วย PyThaiNLP ร่วมกับกฎคำสำคัญ (Keyword-based Rules) และโมเดล SVM ที่ได้รับการฝึกไว้แล้ว
    \item \textbf{การแสดงผลแบบเรียลไทม์:} ระบบแสดงรายการข่าวบนเว็บแอปพลิเคชันพร้อมสรุปเนื้อหา ภาพประกอบ และหมวดหมู่ และแจ้งเตือนข่าวใหม่ทันทีผ่าน WebSocket โดยไม่ต้องรีเฟรชหน้าเว็บ
    \item \textbf{การจัดเก็บข้อมูล:} ระบบจัดเก็บข้อมูลข่าวที่รวบรวมได้ในรูปแบบ JSON บนเครื่องเซิร์ฟเวอร์
    \item \textbf{การตั้งค่าและความยืดหยุ่น:} ระบบรองรับการเพิ่มแหล่งข่าวใหม่ได้โดยไม่ต้องแก้ไขโครงสร้างหลัก ผ่านการลงทะเบียน Source ด้วย Decorator และกำหนดค่าต่าง ๆ ได้ผ่านไฟล์กำหนดค่า เช่น ช่วงเวลาในการดึงข่าว จำนวนข่าวต่อแหล่ง และจำนวนประโยคที่ใช้สรุป
    \item \textbf{ระบบแคช (Cache):} ระบบมีกลไกจัดเก็บข้อมูลข่าวและเนื้อหาที่ดึงมาแล้วไว้ในแคช โดยใช้ตัวระบุเฉพาะของข่าว เช่น URL เป็นกุญแจแคช เพื่อหลีกเลี่ยงการดึงข้อมูลซ้ำจากเว็บไซต์ต้นทางในรอบการทำงานถัดไป ซึ่งช่วยลดภาระของเซิร์ฟเวอร์และเพิ่มความรวดเร็วในการแสดงผล
    \item \textbf{ระบบคุกกี้ (Cookies):} ระบบจัดการและจัดเก็บคุกกี้ของเบราว์เซอร์ร่วมกับ Playwright Service เพื่อคงสถานะเซสชัน (Session) ในการดึงข้อมูลจากเว็บไซต์ข่าวที่แสดงผลเนื้อหาผ่าน JavaScript หรือมีกลไกตรวจสอบเบราว์เซอร์ เช่น Cloudflare ให้สามารถดึงข่าวได้อย่างถูกต้องและต่อเนื่อง
    \item \textbf{คุกกี้เพื่อการปรับแต่งเฉพาะบุคคล (Personalize Cookies):} ระบบใช้คุกกี้ในการจัดเก็บค่ากำหนดของผู้ใช้ เช่น หมวดหมู่ข่าวที่สนใจ ภาษาที่ต้องการแสดงผล และการตั้งค่าส่วนติดต่อผู้ใช้ เพื่อให้ผู้ใช้ได้รับประสบการณ์ที่ปรับแต่งตามความต้องการของแต่ละบุคคล
    \item \textbf{คุกกี้เพื่ออัลกอริทึม (Algorithm Cookies):} ระบบใช้คุกกี้ในการบันทึกพฤติกรรมการอ่านข่าวของผู้ใช้ เช่น หมวดหมู่ที่เข้าชมบ่อย ข่าวที่ใช้เวลาอ่านนาน หรือข่าวที่ผู้ใช้กดบันทึกหรือแชร์ ข้อมูลพฤติกรรมเหล่านี้จะถูกส่งให้อัลกอริทึมแนะนำข่าว (Recommendation Algorithm) เพื่อวิเคราะห์ความสนใจและคัดสรรข่าวที่เกี่ยวข้องมากที่สุดมาแสดงเป็นลำดับแรกบนหน้าเว็บ ซึ่งช่วยเพิ่มความเกี่ยวข้องของเนื้อหาและประสบการณ์การใช้งานที่ดียิ่งขึ้นสำหรับผู้ใช้แต่ละคน
\end{enumerate}

\subsection{ข้อจำกัดของโครงงาน}

\begin{enumerate}[label=\arabic*)]
    \item ระบบรองรับการสรุปและแสดงผลข่าวในภาษาไทยและภาษาอังกฤษเท่านั้น ตามภาษาเดิมของเนื้อหาข่าวแต่ละฉบับ
    \item ระบบครอบคลุมแหล่งข่าวออนไลน์ที่กำหนด และโครงสร้าง HTML หรือ RSS Feed ของเว็บไซต์แหล่งข่าวอาจเปลี่ยนแปลงได้ ทำให้ระบบต้องได้รับการปรับปรุงให้สอดคล้อง
    \item การสรุปข่าวต้องอาศัยการเชื่อมต่ออินเทอร์เน็ตและ API ของโมเดลภาษาขนาดใหญ่ (KKU Gen.AI) หากบริการขัดข้องหรือหมดโควตาการใช้งาน ระบบจะไม่สามารถสรุปข่าวใหม่ได้
    \item ด้วยข้อจำกัดด้านระยะเวลาและงบประมาณ ระบบจึงถูกติดตั้งและทดสอบบนเครื่องคอมพิวเตอร์ภายในมหาวิทยาลัยเท่านั้น ยังไม่ครอบคลุมการใช้งานแบบผู้ใช้พร้อมกันจำนวนมาก
    \item รอบแรกของการดึงข่าวอาจใช้เวลา 2--5 นาที เนื่องจากต้องดึงเนื้อหาบทความของแต่ละข่าวเพื่อใช้ในการสรุป
\end{enumerate}

% -----------------------------------------------------------
\section{ประโยชน์ที่คาดว่าจะได้รับ}
% -----------------------------------------------------------

\begin{enumerate}[label=\arabic*)]
    \item ได้ระบบเว็บแอปพลิเคชันจำแนกและวิเคราะห์ข่าวสารอัตโนมัติแบบเรียลไทม์ ที่ช่วยให้ผู้ใช้ติดตามข่าวสารจากหลายแหล่งได้ในที่เดียวอย่างสะดวก รวดเร็ว และเหมาะสมกับความสนใจของแต่ละบุคคล
    \item ได้แนวทางในการประยุกต์ใช้เทคนิคเว็บสแครปปิง การอ่าน RSS Feed โมเดลภาษาขนาดใหญ่ อัลกอริทึมการกำหนดคุกกี้เพื่อการปรับแต่งเฉพาะบุคคล และอัลกอริทึมแนะนำข่าว เพื่อสร้างระบบวิเคราะห์ข้อมูลข่าวอัตโนมัติ ซึ่งสามารถนำไปพัฒนาต่อยอดกับข้อมูลประเภทอื่น
    \item ได้องค์ความรู้ด้านการประมวลผลภาษาธรรมชาติ การจำแนกข้อความด้วยกฎคำสำคัญร่วมกับโมเดลแมชชีนเลิร์นนิง เทคนิคการสรุปข้อความด้วยโมเดลภาษาขนาดใหญ่ และการพัฒนาระบบเว็บแบบเรียลไทม์
\end{enumerate}
